\section*{Capítulo \ref{ch:b3:computational-chemistry} --- Química computacional: Hartree--Fock y DFT}

\begin{solution}{exo:b3:computational-chemistry:1}
$\dd E/\dd\alpha = \frac32 - \sqrt{2/\pi\alpha} = 0$ da $\alpha = 8/9\pi = 0.283\,
a_0^{-2}$ y $E = \frac32\alpha - 2\sqrt{2\alpha/\pi} = -4/3\pi = \qty{-0.4244}{\hartree}$,
un 15\,\% por encima del valor exacto $-0.5$: una sola gaussiana no tiene ni el punto anguloso ni la
cola de la función $1s$.
\end{solution}

\begin{solution}{exo:b3:computational-chemistry:2}
STO-3G: 5 funciones por C ($1s$, $2s$, tres $2p$) y 1 por H: $6 \times 5 + 6 =
36$. 6-31G(d): por C, 1 (núcleo) + $2 \times 4$ (valencia desdoblada) + 6 ($d$) $= 15$;
por H, 2: $6 \times 15 + 6 \times 2 = 102$.
\end{solution}

\begin{solution}{exo:b3:computational-chemistry:3}
El hamiltoniano electrónico no contiene las masas nucleares: dentro de la
\omterm{def:b3:computational-chemistry:born-oppenheimer}{aproximación de Born--Oppenheimer}, la SEP, y por tanto $r_e$ y $k = U''(r_e)$, es la
misma. $\tilde\omega = \sqrt{k/\mu}/2\pi c$ depende de la \omterm{def:b3:quantum-model-systems:oscillator}{masa reducida}, que se
duplica: razón $\sqrt2 = 1.414$ (medida: 1.413).
\end{solution}

\begin{solution}{exo:b3:computational-chemistry:4}
Ocho átomos dan $3N - 6 = 18$ vibraciones: 17 reales y una imaginaria. Un
\omterm{def:b3:quantum-model-systems:operator}{valor propio} negativo de la hessiana: un punto de silla de primer orden, una estructura
de transición (aquí, por ejemplo, de una rotación o de una eliminación).
\end{solution}

\begin{solution}{exo:b3:computational-chemistry:5}
$\zeta = 27/16$, $E = -(27/16)^2 = \qty{-2.8477}{\hartree} = \qty{-77.49}{eV}$.
Error: $-2.8477 + 2.9034 = \qty{0.0557}{\hartree} = \qty{1.52}{eV}$, es decir,
un 1.9\,\% de la energía total ---pero más que muchas energías de reacción.
\end{solution}

\begin{solution}{exo:b3:computational-chemistry:6}
$(-1 - E)(-0.5 - E) - (-0.2 - 0.3E)^2 = 0$, es decir, $0.91E^2 + 1.38E + 0.46 = 0$:
$E = \qty{-1.0217}{\hartree}$ y \qty{-0.4947}{\hartree}. Sí: incorporar la
segunda función baja la energía por debajo de $H_{11}$, como permite el \omterm{thm:b3:computational-chemistry:variational}{principio
variacional}.
\end{solution}

\begin{solution}{exo:b3:computational-chemistry:7}
$(102/36)^4 = 64$: unos 640 minutos, unas 11 horas.
\end{solution}

\begin{solution}{exo:b3:computational-chemistry:8}
$0.5782 \times 27.211 = \qty{15.73}{eV}$, 0.30 eV por encima del valor medido. Orbitales
congelados: el ion real se relaja y queda más bajo, de modo que el valor de Koopmans es demasiado alto.
Correlación ausente: la molécula neutra (dos electrones) tiene más \omterm{def:b3:computational-chemistry:correlation}{energía de
correlación} que el ion (un electrón), lo que hace mayor el valor verdadero.
Aquí gana el primer error.
\end{solution}

\begin{solution}{exo:b3:computational-chemistry:9}
$-0.5960 - 2(-0.4666) = \qty{0.3372}{\hartree} = \qty{9.18}{eV}$ por encima de dos átomos.
La función RHF conserva un 50\,\% de \ce{H+ H-}, y separar un protón y un
hidruro cuesta la diferencia entre la energía de ionización de H y la
afinidad electrónica de H, que la energía incorpora en la media.
\end{solution}

\begin{solution}{exo:b3:computational-chemistry:10}
$E(c_1^2 + 2c_1c_2S + c_2^2) = c_1^2H_{11} + 2c_1c_2H_{12} + c_2^2H_{22}$.
Derivando respecto de $c_1$ con $\partial E/\partial c_1 = 0$:
$E(2c_1 + 2c_2S) = 2c_1H_{11} + 2c_2H_{12}$, es decir, $(H_{11} - E)c_1 + (H_{12} -
ES)c_2 = 0$; del mismo modo, $(H_{12} - ES)c_1 + (H_{22} - E)c_2 = 0$.
\end{solution}

\begin{solution}{exo:b3:computational-chemistry:11}
$k \approx [E(1.30) - 2E(1.35) + E(1.40)]/(0.05)^2 = 0.001417/0.0025 =
\qty{0.567}{\hartree\per\bohr\squared} = \qty{882}{N/m}$. Con $\mu = m_H/2 =
\qty{8.37e-28}{kg}$, $\omega = \sqrt{k/\mu} = \qty{1.03e15}{s^{-1}}$ y
$\tilde\omega = \omega/2\pi c = \qty{5452}{cm^{-1}}$, un 24\,\% por encima del
$\tilde\omega_e$ medido: una curva RHF en \omterm{def:b3:computational-chemistry:basis}{base mínima} es demasiado empinada (sin correlación, con una base
demasiado pequeña), y por eso las frecuencias calculadas se escalan a la baja.
\end{solution}

\begin{solution}{exo:b3:computational-chemistry:12}
Unas diferencias de pocos \unit{kJ/mol} en las que interviene un enlace de hidrógeno exigen correlación y
dispersión: un funcional híbrido con corrección de dispersión, o mejor un
método de función de onda correlacionado para las energías finales, con una base
triple desdoblada polarizada que incluya \omterm{def:b3:computational-chemistry:split-valence}{funciones difusas}. Se optimizan ambos confórmeros al
mismo nivel, se confirman los mínimos con las frecuencias (todas reales), se añaden las \omterm{def:b3:quantum-model-systems:oscillator}{energías de punto
cero}, se pone a prueba la base con un cálculo mayor y se compara con un
sistema emparentado de energía conformacional conocida.
\end{solution}

\begin{solution}{pb:b3:computational-chemistry:1}
\textbf{1.} Cada \omterm{def:b3:computational-chemistry:basis}{orbital de tipo Slater} de una \omterm{def:b3:computational-chemistry:basis}{base mínima} se sustituye por una
contracción fija de tres gaussianas ajustadas a él.
\textbf{2.} El núcleo de helio atrae más a sus electrones, de modo que su función $1s$
es más compacta (exponentes mayores). $6.3624/3.4253 = 1.857 = (1.69/1.24)^2$.
\textbf{3.} Dos funciones de base, dos electrones, un orbital ocupado (y uno
virtual).
\textbf{4.} $V_{\mathrm{nn}} = Z_{\mathrm{He}}Z_{\mathrm H}/R = 2/1.4632 =
\qty{1.3669}{\hartree}$.
\textbf{5.} Las dos funciones solapan mucho (54\,\%): los átomos están lo bastante cerca
para enlazarse.
\textbf{6.} $h_{11}$ es la energía de un electrón en la función de He con los dos
núcleos pero sin ningún otro electrón; el núcleo de He (carga 2) lo retiene con mucha más
fuerza.
\textbf{7.} $(h_{11} - \varepsilon)(h_{22} - \varepsilon) - (h_{12} - \varepsilon S)^2 = 0$:
$0.71185\varepsilon^2 + 2.79292\varepsilon + 2.44969 = 0$.
\textbf{8.} $\varepsilon = \qty{-2.600}{\hartree}$ y \qty{-1.324}{\hartree}.
\textbf{9.} $\mathbf h$ ignora la repulsión entre los dos electrones; el orbital de
núcleo es demasiado contraído y demasiado bajo.
\textbf{10.} $F_{\mu\nu} = h_{\mu\nu} + \sum_{\lambda\sigma}P_{\lambda\sigma}[(\mu\nu|\sigma\lambda) -
\frac12(\mu\lambda|\sigma\nu)]$, con $P_{\lambda\sigma} = 2C_{\lambda1}C_{\sigma1}$.
\textbf{11.} La repulsión coulombiana de cada electrón por la nube de carga del
otro (y, en general, el intercambio).
\textbf{12.} $\mathbf F$ depende de $\mathbf P$, que depende de los orbitales obtenidos
a partir de $\mathbf F$: las ecuaciones no son lineales y se resuelven por aproximaciones
sucesivas hasta la autoconsistencia.
\textbf{13.} $c_1^2 + c_2^2 + 2c_1c_2S = 0.7684 + 0.0410 + 0.1906 = 1.0000$.
\textbf{14.} He: $1.5369 + 0.1906 = 1.727$; H: $0.0820 + 0.1906 = 0.273$
electrones.
\textbf{15.} Cargas: He $+0.27$, H $+0.73$. La carga positiva está sobre todo en el
hidrógeno: \ce{HeH+} es un protón unido a un átomo de helio, \ce{He + H+}, como cabía esperar,
ya que la energía de ionización de He (24.6 eV) supera con mucho la de H (13.6 eV).
\textbf{16.} $-\varepsilon_1 = \qty{1.633}{\hartree} = \qty{44.4}{eV}$.
\textbf{17.} $E_{\mathrm{el}} = -2.8418 - 1.3669 = \qty{-4.2087}{\hartree}$.
\textbf{18.} Tres ciclos alcanzan \qty{-2.8418}{\hartree}. Cada ciclo da un
determinante, cuya energía es una cota superior de la energía de Hartree--Fock
convergida (\omterm{thm:b3:computational-chemistry:variational}{principio variacional}), y las iteraciones la mejoran.
\textbf{19.} La base con $\zeta = 2.0925$ da la energía más baja, luego es la
mejor para esta molécula (\omterm{thm:b3:computational-chemistry:variational}{principio variacional}): la función del helio
es más compacta en \ce{HeH+} que en el átomo libre, cuyo $\zeta = 1.69$ copia la
base estándar.
\textbf{20.} El orbital antienlazante vacío $\sigma^*$; su energía aproximaría
menos la afinidad electrónica de \ce{HeH+} (mal, en una base tan pequeña).
\textbf{21.} Cero (ningún electrón); $\qty{-2.8078}{\hartree}$.
\textbf{22.} $-2.8078 - (-2.8418) = \qty{0.0340}{\hartree} = \qty{89}{kJ/mol}$.
\textbf{23.} Alrededor de la mitad del valor medido: una \omterm{def:b3:computational-chemistry:basis}{base mínima} no puede polarizar la
densidad del helio hacia el protón (no hay funciones $p$), de modo que el enlace
es demasiado débil. (Las correcciones de punto cero y a \qty{298}{K} son pequeñas a su lado.)
\textbf{24.} $-(24.5874 + 54.4178)/27.211 = \qty{-2.9034}{\hartree}$; el átomo STO-3G
queda \qty{0.0956}{\hartree} (\qty{2.6}{eV}) demasiado alto.
\textbf{25.} Las geometrías dependen de cómo cambia la energía con $R$; la mayor parte del
error, concentrada en el núcleo, es casi la misma en todas las geometrías y
se cancela.
\textbf{26.} \textbf{$E(\ce{HeH+}) = \qty{-2.8418}{\hartree}$ en RHF/STO-3G (base
estándar), a \qty{1.4632}{\bohr}} (y \qty{-2.8607}{\hartree} con el clásico
$\zeta_{\mathrm{He}} = 2.0925$).
\end{solution}
